\chapter{Working at Height Risk Assessment}
\label{chap:Working at Height Risk Assessment}

%---------------------------- SECTION ----------------------------
\section{Why this assessment matters in construction}

\paragraph{}
Falls from height are the leading cause of fatal injury in the United Kingdom construction sector and have remained so for every year in which the HSE has published disaggregated fatality statistics. Approximately one-third of all construction fatalities in any given year result from a fall --- a proportion that has remained remarkably persistent despite decades of regulatory development, industry guidance, and enforcement activity. The work at height risk assessment is therefore not merely the most frequently required risk assessment on a construction site; it is the assessment whose quality and rigour are most directly correlated with the likelihood of a worker surviving the project. The Work at Height Regulations 2005 (WAHR 2005) impose a duty on every employer and self-employed person to take all reasonable steps to prevent any person falling a distance liable to cause personal injury, and the regulations apply to any work at height regardless of the duration of the activity or the height involved. A tradesperson standing on a stepladder to reach a high-level fixing, a scaffold erector boarding out at 15\,m, a roofer stripping slates on a 30-degree pitch, and a steeplejack carrying out inspection of a heritage chimney stack are all working at height and all require individual, competent risk assessment.

\paragraph{}
The regulatory and practical landscape for working at height on construction sites is more complex and more extensively codified than for many other hazard categories. This complexity arises because the range of access methods available to construction workers --- traditional tube-and-fitting scaffold, proprietary system scaffold, mobile elevated work platforms (MEWPs), mobile access towers, ladders, rope access, fixed and portable work platforms, roof edge protection systems, and personal fall protection systems --- is wide, and because each method carries its own specific hazard profile, competence requirements, inspection regime, and regulatory basis. No single document addresses all these methods: the NASC TG20:21 guide addresses tube-and-fitting scaffold; NASC SG4:22 addresses prevention of falls in scaffolding operations; PASMA codes of practice address mobile towers; IPAF guidance addresses MEWPs; BS~EN~131 and the Ladder Association's guidance address ladders; and the comprehensive suite of EN~363 fall protection system standards addresses personal fall protection. The working at height risk assessor must have sufficient knowledge across this entire landscape to select the most appropriate method for each task and to specify the controls applicable to the chosen method with the required level of technical accuracy.

\paragraph{}
The interface between working at height and other hazard categories creates compounding risks that the assessment must specifically address. Operatives working at height are simultaneously exposed to the physical risk of the fall and to the task-specific hazards of the activity they are carrying out at height: a welder on a scaffold is exposed to both the risk of falling and the risk of arc flash, UV exposure, and fire from hot works; a mechanical and electrical operative on a MEWP is exposed to both the risk of falling from the platform and the risk of a collision between the machine and other plant; a roofer on a fragile surface is exposed to the risk of both a fall through the surface and a fall from the edge. The assessment must address each of these interactions rather than treating the working at height hazard in isolation from the trade activities being performed. The relationship between working at height and lifting operations is particularly significant for MEWPs, which are covered in detail in Chapter~\ref{chap:Lifting Operations and Lifting Equipment Risk Assessment} and are cross-referenced here where MEWPs are used as a working at height access method.

\paragraph{}
The legal, technical, and human factors aspects of working at height risk are mutually dependent. The best-designed scaffold, equipped with compliant edge protection and inspected to the required frequency, will not prevent a fall if the worker using it does not understand the importance of maintaining three points of contact on the access ladder, does not check the scaffold tag before using the structure, or removes a guardrail temporarily to improve access and fails to reinstate it. Conversely, the most motivated and safety-aware worker cannot compensate for a structurally defective scaffold, an uninspected MEWP, or an anchor point that is not capable of sustaining the loading imposed by a fall arrest event. The working at height risk assessment must therefore address both the technical specification of the access system and the human factors that govern whether specified controls will be implemented consistently in practice.

\barquote{``Every fall from height is preventable. The challenge is not technical knowledge --- it is the discipline to apply that knowledge without exception, on every job, every day, including the last job of a long Friday shift.'' --- NASC, \textit{Safety Guidance Note SG4:22 --- Preventing Falls in Scaffolding Operations}, 2022}

\example{Working at Height Incident Scenario}{A groundworks subcontractor engaged in the construction of a large logistics warehouse was installing ground beams adjacent to the building's perimeter. A temporary access arrangement had been constructed using scaffold boards laid across scaffold tubes supported on adjustable bases --- a arrangement that had not been designed by a competent person and had not been formally inspected. A groundworks operative used this arrangement to access a position approximately 1.8 metres above the surrounding ground to mark out anchor bolt positions. The boards had not been secured to the support tubes and, under the operative's weight combined with a lateral movement as he reached for a fixing, the boards displaced. The operative fell approximately 1.8\,m, striking the edge of a concrete ground beam with his shoulder and head. He sustained a serious head injury requiring hospitalisation and suffered permanent partial hearing loss. Investigation found: no risk assessment for the temporary access arrangement; no competent person involvement in its design; no inspection before use; no edge protection to prevent fall from the platform level; and no formal working at height permit or method statement covering the activity. The task was assessed internally as ``not really height work'' because the height was under 2 metres, reflecting the common and legally incorrect belief that the Work at Height Regulations 2005 only apply above a threshold height.}

%---------------------------- SECTION ----------------------------
\section{Legal and professional framework}

The statutory framework for working at height in Great Britain is anchored by the Work at Height Regulations 2005, which implement the European Temporary or Mobile Construction Sites Directive and the Work Equipment Directive in the specific context of work at height. The regulations apply to all work at height where there is a risk of a fall liable to cause personal injury, with no minimum height threshold. The employer and self-employed person must comply with the requirements of the regulations for all working at height activities, regardless of duration. The WAHR 2005 established the regulatory hierarchy --- avoid, prevent falls, mitigate the consequences of falls --- which provides the evaluative framework for selecting access methods and protection systems and is reflected in the hierarchy of controls set out in Section~\ref{sec:wah_hierarchy} of this chapter.

The professional and technical framework is supported by an extensive body of standards, codes of practice, and industry guidance that translates the regulatory requirements into specific technical criteria for the design, erection, inspection, and use of access equipment:

\begin{itemize}
  \bitem{Work at Height Regulations 2005 (WAHR 2005)}{The principal statutory instrument governing working at height across all sectors in Great Britain. Regulation~6 requires employers to avoid work at height where reasonably practicable; where it cannot be avoided, to use work equipment or other measures to prevent falls; and where falls cannot be prevented, to use work equipment or other measures to minimise the distance and consequences of any fall that does occur. Regulation~13 and Schedule~7 set out specific requirements for the selection of work equipment for use at height, and Schedule~2 sets out requirements for the inspection of work equipment used at height and the recording of inspection results.}

  \bitem{TG20:21 --- Guide to Good Practice for Scaffolding with Tubes and Fittings (NASC)}{The National Access and Scaffolding Confederation's comprehensive technical guide to the design and erection of tube-and-fitting scaffold for use on construction sites. TG20:21 is the definitive UK industry reference for compliant scaffold design and provides standard configurations for a wide range of scaffolding scenarios. TG20:21 compliance provides a defence against allegations of inadequate scaffold design and is widely specified in principal contractors' supply chain requirements. The 2021 edition incorporates updated guidance on system scaffolding, edge protection standards, and the advanced guardrail (AGR) technique.}

  \bitem{SG4:22 --- Preventing Falls in Scaffolding Operations (NASC)}{The NASC Safety Guidance Note specifically addressing the risk of falls during the erection, alteration, and dismantling of scaffolding. SG4:22 establishes the requirement for the advanced guardrail (AGR) technique as the primary method for preventing falls during scaffold boarding-out operations, and specifies the sequence of operations, equipment, and supervision required for compliant safe-system scaffold erection and dismantling. It also addresses the use of personal fall protection systems by scaffold erectors as a supplementary measure where the AGR technique cannot be fully applied.}

  \bitem{PASMA Code of Practice for Mobile Access Towers}{The Prefabricated Access Suppliers' and Manufacturers' Association's technical and training standards for the assembly, use, and dismantling of mobile aluminium alloy access towers. PASMA training is recognised as the minimum competence standard for operatives assembling and using mobile towers, and the PASMA Operator Card provides evidence of training completion. The PASMA code specifies maximum working heights, stability criteria, outrigger requirements, castors and locking arrangements, and inspection requirements for mobile towers.}

  \bitem{BS EN~131:2015+A1:2019 --- Ladders (BSI)}{The current European standard for portable ladders, specifying design requirements, load classifications, test methods, and marking requirements. BS EN~131 replaced the earlier BS~2037 (Class~1 Industrial) and BS~2037 (Class~3 Domestic) series for new ladders. The standard distinguishes between professional-duty (Class~1) ladders with a maximum user weight of 150\,kg and non-professional (Class~EN) ladders not appropriate for construction use. Ladder selection, inspection, and use on construction sites must comply with BS EN~131 requirements and the Ladder Association's guidance on safe ladder use.}

  \bitem{BS EN~363:2008 --- Personal Fall Protection Systems (BSI)}{The European standard defining the system components and requirements for personal fall protection, including fall arrest systems, work restraint systems, work positioning systems, rope access systems, and rescue systems. BS EN~363 provides the framework within which the individual component standards --- EN~361 (harnesses), EN~355 (energy-absorbing lanyards), EN~360 (retractable type fall arresters/inertia reels), EN~795 (anchor devices), EN~354 (lanyards) --- sit. The standard establishes compatibility and system requirements that must be addressed in the selection and specification of personal fall protection equipment.}

  \bitem{BS EN~795:2012 --- Personal Fall Protection Equipment --- Anchor Devices (BSI)}{Specifies design and performance requirements for anchor devices used in personal fall protection systems. Categorises anchor devices into Classes~A through E based on structural attachment, which is critical for determining appropriate anchor points on construction structures. Class~A structural anchors (permanently fixed to walls, roofs, or other rigid structures) must be designed and installed by a competent person with knowledge of the loading imposed by fall arrest events (typically 6\,kN minimum), and must be inspected by a competent person annually and after any fall event.}

  \bitem{BS EN~1263:2014 --- Safety Nets --- Safety Requirements and Test Methods (BSI)}{Specifies requirements for safety nets used as collective protection measures to arrest falls of persons and materials on construction sites. Part~1 covers safety requirements and test methods; Part~2 covers erection requirements. Safety nets conforming to BS EN~1263 provide a collective fall mitigation measure that does not depend upon individual workers wearing or connecting to personal fall protection equipment. The standard specifies maximum mesh size, energy absorption, and installation clearance requirements that must be verified by a competent net erector.}
\end{itemize}

%---------------------------- SECTION ----------------------------
\section{Conducting the assessment using the five-step method}

\subsection{Step 1 --- Identify Hazards}

\paragraph{}
The hazard identification phase for a working at height risk assessment requires a systematic examination of all locations on the site where work will be carried out at height or where access to height is required, together with all the access methods and work activities involved. A site-wide working at height register should catalogue every area where height work will occur: facade and structural scaffold at perimeter elevations; internal scaffold within the structural frame; floor edge conditions during concrete pouring and structural steel erection; roof work on permanent roofing, cladding, and waterproofing systems; maintenance of temporary works at height; loading and unloading from elevated platforms; cleaning and inspection activities on completed sections; and access to the tops of temporary works such as cranes, hoists, and mast climbers. For each identified location and activity, the specific hazard must be characterised: an unprotected floor edge presents a different hazard profile from a fragile roof surface, which differs from the risk of falling through an open shaft, which differs from the risk of falling from a ladder during a brief access task. The assessment must also identify the dynamic changes to the site's working at height hazard profile that will occur as the project progresses --- floor edges that are protected as the structure reaches each level, scaffolding that is progressively erected and modified, and access routes that change as the building progresses through phases.

\subsection{Step 2 --- Decide Who or What Is Affected}

\paragraph{}
The population at risk from falls at height on a construction site encompasses every person who sets foot on the site, since unprotected edges, open voids, and unguarded excavations may exist in locations that are traversed during normal site access as well as in the work areas themselves. Scaffold erectors (scaffolders) are at highest risk during the erection, alteration, and dismantling process --- the period during which the protection system that will ultimately protect all other workers does not yet exist in its final form. Structural steelworkers working on the advancing steel frame before floors and edge protection are installed face the highest inherent height risk of any trade on a construction project. Roofers, cladding operatives, and waterproofing contractors work at the most exposed heights and on surfaces that may be fragile, wet, or steeply inclined. General operatives, labourers, and tradespeople carrying out work adjacent to unprotected edges --- even briefly, to pass materials or check measurements --- face the cumulative probability of a lapse in attention or an unexpected disturbance leading to a fall. Supervisory and management staff visiting work areas to inspect progress must be included in the assessment: a site manager who routinely accesses areas without appropriate edge protection in place is both personally at risk and implicitly normalising non-compliant behaviour. Visiting professionals, including architects, structural engineers, and client representatives, must be included in site-specific induction that covers working at height restrictions in their specific visiting areas.

\subsection{Step 3 --- Evaluate Likelihood and Consequence}

\paragraph{}
The evaluation of fall risk must reflect the height from which a fall could occur, the nature of the landing surface, the posture and activity of the worker at the moment of a potential fall, and the effectiveness of collective and individual protection in place. Falls from any height carry a realistic risk of serious injury or death, and the HSE's data demonstrate that fatal falls in construction occur from heights as low as 2\,metres. The consequence dimension of the risk matrix should therefore be assessed as High or Very High for all unprotected working at height scenarios regardless of the specific height, since even a fall from standing height onto a hard surface or onto projecting reinforcement or other sharp hazards can be fatal. Likelihood must be assessed against the specific site conditions: open floor edges without guardrails during a concrete pour are a high-probability hazard; an inspected scaffold with compliant edge protection in good condition is a substantially reduced-probability hazard but not zero risk, since falls can occur from scaffold through distraction, carrying of loads, or loss of footing on a wet boarding surface. The assessment must incorporate the specific likelihood factors of the construction phase: structural steel erection has an inherently higher probability of a fall than the fit-out phase of the same building, because structural steel working is inherently more exposed and because the protective systems provided by the building's structure are not yet in place.

\subsection{Step 4 --- Select and Implement Controls}

\paragraph{}
Control selection for working at height must follow the WAHR~2005 regulatory hierarchy (avoid, prevent, mitigate) as the primary evaluative framework, selecting the highest-order control that is technically and practically achievable for the task. Collective measures that protect all workers in the area --- edge protection, safety nets, scaffolding platforms --- must be preferred over individual personal fall protection systems that protect only the wearer and depend upon correct selection, fitting, connection, and maintenance for their effectiveness. Where scaffolding is selected as the primary access and edge protection method, it must be designed and erected in accordance with TG20:21, with edge protection complying with the WAHR~2005 Schedule~2 requirements: a top guardrail at a minimum height of 950\,mm, an intermediate guardrail (not more than 470\,mm below the top guardrail), and a toe board of minimum 150\,mm height. Personal fall protection systems, where specified, must be selected from a compatible system in accordance with BS EN~363, with anchor points designed to BS EN~795, and operatives trained and certificated to use the system. Mobile access towers must be assembled and used in accordance with PASMA standards by PASMA-trained operatives. Ladders should be restricted to access and egress tasks of short duration and must not be used as work platforms for sustained work activities. Where fragile roofs or fragile rooflights are present, the guidance in the WAHR~2005 Schedule~2 regarding warning notices, barriers, and crawling boards must be applied, and the fuller treatment of roofing and fragile surface risk is addressed in Chapter~\ref{chap:Roofing and Fragile Surfaces Risk Assessment}.

\subsection{Step 5 --- Record and Review}

\paragraph{}
The written working at height risk assessment must be specific to the task, trade, and location --- a generic ``working at height'' assessment that does not differentiate between scaffold erection, MEWP use, roof work, and ladder access tasks does not constitute a suitable and sufficient assessment under the WAHR~2005. The assessment must document the selected access method, the specific edge protection or fall arrest system specified, the inspection requirements and frequency, the competence requirements for the operatives involved, and the supervision arrangements. Scaffold inspection records must be maintained in the format required by WAHR~2005 Schedule~7 and TG20:21: formal inspection of completed scaffold before first use; repeated inspection at intervals not exceeding seven days; and inspection following any event likely to have affected the scaffold's stability or integrity, including adverse weather, plant impact, and significant loading. The inspection records must be signed by a competent person and retained for three months after the scaffold has been dismantled. MEWP pre-use inspections must be documented in the machine's daily inspection log. Personal fall protection equipment must be recorded in an individual equipment log, with evidence of manufacturer's inspection intervals, visual pre-use checks, and any post-fall inspection by a competent person.

\example{Five-Step Application}{A principal contractor constructing a 12-storey reinforced concrete residential tower carries out the working at height risk assessment as follows. \textbf{Step~1 (Hazards):} Open floor edges at each level during concrete pours (levels 1--12 progressing at 2-week cycle); open core and lift shaft voids throughout the structure; perimeter scaffold for facade cladding (aluminium composite panel) from ground to full height; internal scaffold for fit-out work in common areas; roofing works on flat roof with bituminous waterproofing at level~12; MEWPs for external cladding fixings and sealant application; temporary works at height for prop and shutter erection; steel fixing and concrete pour activities in confined reinforced concrete core; access ladders for supervisory inspection rounds. \textbf{Step~2 (Who Affected):} Structural concrete gangs (formwork, rebar, pour); perimeter scaffold erectors (NASC-registered subcontractor); cladding operatives; internal fit-out trades at all levels; building services engineers on MEWPs; visiting structural engineers and architect's representatives; principal contractor's site management team. \textbf{Step~3 (Evaluation):} Open floor edges during pour --- Very High inherent risk; open core void before permanent protection --- Very High; MEWP at perimeter --- High inherent; ladder access during inspection --- Medium. \textbf{Step~4 (Controls):} Proprietary safety decking system with integral edge protection installed before each floor pour commences; core void protection using proprietary hop-up platforms within the core; TG20:21-compliant perimeter scaffold with standard guardrail and toe board; MEWP selection from IPAF-registered operator only (IPAF~3a and~3b cards); all MEWPs to be equipped with secondary guarding on boom; short-duration ladder access only for inspection rounds with competent person designation; EN~361 harness and EN~795 Class~A roof anchor for roofers on level~12 perimeter. \textbf{Step~5 (Record):} Working at height assessment for each trade package signed by subcontractor supervisor and principal contractor's site manager; scaffold inspection register maintained at 7-day intervals; MEWP daily pre-use log maintained by operator; level-by-level floor edge protection sign-off record before each structural concrete pour commences. Residual risk for floor edges rated Low with proprietary decking system; open core void rated Medium pending permanent protection installation.}

%---------------------------- SECTION ----------------------------
\section{Applying the hierarchy of controls}
\label{sec:wah_hierarchy}

\paragraph{}
The hierarchy of controls for working at height is established directly by the Work at Height Regulations 2005 in Regulation~6, which creates a three-level hierarchy: first, avoid work at height where reasonably practicable; second, where height work cannot be avoided, prevent falls using suitable work equipment or other measures; third, where falls cannot be entirely prevented, minimise the distance and consequences of a fall that does occur. This regulatory hierarchy directly maps onto the general hierarchy of controls framework and must be the primary evaluative structure used in the assessment. Controls must not be selected on the basis of cost, convenience, or familiarity alone --- the assessment must demonstrate that higher-order controls were genuinely considered and either implemented or rejected with clear documented justification.

\begin{itemize}
  \bitem{Elimination}{The most effective control is to eliminate the need to work at height entirely, by redesigning the task, the sequence of works, or the building design so that the activity can be carried out at ground level. Pre-fabricated and pre-assembled components --- structural steel connections assembled at ground level before crane lifting; rooflight frames glazed at ground level before installation; pre-wired mechanical and electrical modules assembled in workshops before hoisting into position --- reduce the extent of height work on site and should be considered during the CDM design stage. Where elimination is not possible for the overall activity, it may be possible to eliminate the height work element of a specific sub-task by using extended tools, pole-mounted equipment, or robotic devices to carry out fixings, cleaning, or inspection from ground level.}

  \bitem{Substitution}{Where complete elimination is impracticable, the access method selected should represent the lowest-risk option capable of delivering the required work. A MEWP providing a stable, enclosed working platform is a lower-risk access method for a specific external sealant application task than a ladder --- substitution of the ladder with the MEWP reduces the risk of the working at height element regardless of the additional hazards introduced by the MEWP itself (addressed in Chapter~\ref{chap:Lifting Operations and Lifting Equipment Risk Assessment}). A proprietary system scaffold, providing modular and dimensionally consistent edge protection, may represent a lower-risk substitution for a bespoke tube-and-fitting scaffold in environments where tube-and-fitting erection at the specific geometry is technically complex and the risk of erection error is elevated.}

  \bitem{Engineering controls}{Engineering controls for prevention of falls include: edge protection systems (guardrails, intermediate rails, toe boards) to WAHR~2005 Schedule~2 specification; safety netting to BS EN~1263 installed beneath the working level during structural steel erection and roof decking; proprietary safety decking systems that provide an integral working platform and edge protection during concrete floor construction; rigid covers and barriers over open voids, shafts, and staircase openings maintained in place at all times except during active material or person transit; and fixed access ladder systems with hooped safety protection or integrated fall arrest rails for permanent vertical access. Engineering controls are the preferred protection method because they provide collective protection to all workers in the area without requiring individual action or consistent behaviour.}

  \bitem{Administrative controls}{Administrative controls for working at height include: permit-to-work systems for high-risk height activities such as scaffold erection on live highways or in windy conditions; restricted access zones beneath scaffold erection and dismantling operations; mandatory site induction covering working at height restrictions and edge protection requirements; seven-day scaffold inspection regime with signed records; daily pre-use inspection of MEWPs, mobile access towers, and personal fall protection equipment; restricted use of ladders to brief access tasks with a signed site rule that ladders are not to be used as work platforms; briefing of all trades on the requirement to report defects in edge protection immediately and not to remove or modify guardrails without authorisation; and the SG4:22 advanced guardrail sequence applied as a mandatory method for scaffold erection and dismantling.}

  \bitem{PPE}{Personal fall protection systems, comprising harness, lanyard or inertia reel, and anchor device, are the last line of defence against falls at height and must not be specified as the primary control where collective protective measures are practicable. Where personal fall protection is specified, the full system must be designed to BS EN~363: the harness must be to BS EN~361, the energy-absorbing lanyard to BS EN~355 (or the inertia reel to BS EN~360), the anchor device to BS EN~795 of the appropriate class for the attachment point, and the complete system must provide a fall clearance distance sufficient to ensure the operative cannot strike any lower level or obstruction before the fall is arrested. Operatives must be trained and competent in the correct donning, connection, and pre-use inspection of the system, and the equipment must be individually logged, inspected at manufacturer's recommended intervals, and withdrawn from service immediately following any fall event.}
\end{itemize}

%---------------------------- SECTION ----------------------------
\section{Cadence of assessment and review triggers}

\paragraph{}
The working at height risk assessment is among the most frequently reviewed of all construction risk assessments, because the physical nature of the height hazard changes continuously as the building's structure rises, as new trades access different areas, as scaffold is erected and modified, and as temporary edge protection systems are progressively replaced by permanent structural edge conditions. The WAHR~2005 requires that work equipment used at height, including scaffolding, is inspected before first use by any person, at seven-day intervals, and after any event likely to have affected its stability. This statutory inspection requirement creates a minimum review cadence that must be augmented by the event-triggered reviews described below.

\begin{itemize}
  \bitem{Commencement of a new trade activity at height}{Each new trade mobilisation --- structural steel erectors, cladding operatives, roofers, MEWP operators, scaffold erectors --- introduces different height hazards, different access methods, and different competence requirements. A trade-specific working at height assessment must be reviewed or produced before the trade commences work, and the principal contractor must satisfy itself that the subcontractor's own assessment is adequate for the specific site conditions.}

  \bitem{Structural or scaffold configuration change}{Any modification to the scaffold configuration, including the addition of new lifts, the removal of ties, the introduction of loading bays, or the alteration of the access route, must be followed by a formal scaffold inspection before the modified scaffold is returned to use. The scaffold inspection record must reflect the specific configuration change and confirm that the modified structure complies with TG20:21 or the approved design.}

  \bitem{Adverse weather}{High wind, ice, frost, snow on scaffold boards, and heavy rainfall all materially increase the risk of a fall from height. Following any period of adverse weather that may have affected scaffold stability or working surface conditions, a competent person must inspect the scaffold and all height access equipment before work resumes. Sites operating in winter conditions must have specific adverse weather protocols that specify the wind speed threshold at which MEWP and mobile tower use is prohibited.}

  \bitem{Any fall from height or near miss}{Any fall from height, regardless of whether it results in injury, and any near miss where a fall nearly occurred, must trigger an immediate review of the working at height assessment and the specific control measures in place at the incident location. The review must identify whether the failure was in the design of the system, in the inspection regime, in the training and competence of the operative, or in the supervision arrangements, and must implement corrective actions before work in the affected area resumes.}

  \bitem{Change in working method or programme sequence}{Programme changes that require work at height to commence in areas not yet equipped with the planned protective measures --- for example, trades accessing upper floor areas before the perimeter scaffold has reached that level --- must trigger a working at height assessment review and must not be authorised without an identified compliant alternative access method and edge protection arrangement.}
\end{itemize}

%---------------------------- SECTION ----------------------------
\section{Common failure modes}

\begin{itemize}
  \bitem{Removal and non-reinstatement of edge protection}{Guardrails and intermediate rails on scaffold are routinely removed to facilitate material loading, delivery, or the movement of large items onto the scaffold, and subsequently not reinstated. The protection gap may persist for hours or even days, during which all workers using the scaffold are exposed to an unprotected edge. A site rule requiring that any removal of edge protection must be authorised by the site manager, carried out by a competent scaffold-trained operative, and reinstated immediately upon completion of the loading activity --- not left for a later visit --- is the administrative control for this failure mode.}

  \bitem{Use of ladders as work platforms}{Operatives routinely use ladders as work platforms for sustained tasks --- fixing, painting, pipework installation --- because ladders are immediately available and avoidance of them seems to create delay. The Work at Height Regulations 2005 Schedule~6 specifies that ladders may only be used as work equipment at height where the risk assessment demonstrates that use of more suitable equipment is not justified because of the low risk and short duration of use or existing site features. Ladders used as work platforms for anything other than brief access or a genuinely short-duration, low-risk task are non-compliant. The training, induction, and supervision response must consistently challenge and redirect this behaviour.}

  \bitem{Scaffold inspection records not maintained or not site-specific}{A substantial proportion of scaffold inspection failures identified in HSE enforcement actions involve inspection records that are either absent entirely, completed as carbon-copy identical entries without genuine engagement with the scaffold's condition, or completed by individuals who do not possess the competence level required by WAHR~2005 for scaffold inspection. The inspection record must be specific to the scaffold being inspected, must identify any defects found and the actions taken, and must be signed by a named competent person. A scaffold inspector must be able to demonstrate both theoretical competence (typically through a Scaffold Inspection competence qualification or equivalent) and current practical familiarity with the specific scaffold configuration.}

  \bitem{Personal fall protection not connected or incorrectly connected}{Where personal fall protection is specified and provided, it is routinely observed in use either disconnected (the operative has unclipped the lanyard for convenience), connected to an anchor of inadequate strength (a scaffold tube rather than a BS EN~795 Class~A structural anchor), or connected with insufficient fall clearance distance to arrest the fall before the operative strikes a lower level. The provision of personal fall protection equipment does not constitute effective implementation of the control: it requires training, supervision, and the physical provision of suitable anchor points before the work commences.}

  \bitem{Mobile access towers assembled by untrained operatives}{PASMA statistics and HSE investigation records consistently identify mobile access tower incidents where the tower has been assembled by an operative who has not received PASMA training, resulting in errors such as: failure to use outriggers on the appropriate configuration; incorrect bay dimensions exceeding the manufacturer's safe height-to-base ratio; unlocked or missing castors; missing end frames or cross braces; or incomplete installation of the working platform with through-access hatch. PASMA training should be mandatory for all operatives assembling or using mobile access towers, and the assembled tower should be inspected by a competent person before first use.}

  \bitem{Inadequate clearance distance for fall arrest}{A fall arrest system that, when activated, allows the operative to strike the level below before the fall is arrested has failed as a control despite being technically ``in place.'' The calculation of fall clearance distance --- the sum of the height of the anchor above the operative, the length of the lanyard, the extension of the energy absorber, the height of the operative's centre of gravity, and a safety margin --- must be carried out for every fall arrest application and the result verified against the available clearance. In environments where the available clearance is insufficient for a conventional fall arrest system, a work restraint system (preventing the operative from reaching the fall hazard) or an inertia reel system with a shorter arrest distance may be more appropriate.}
\end{itemize}

\example{Failure Pattern}{The HSE's investigation of a fatal fall from scaffold on a domestic refurbishment project identified the following failure chain. The project involved the replacement of windows on the upper two floors of a terraced Victorian property, with a three-lift tube-and-fitting scaffold erected by a sole trader scaffolder who was a sub-subcontractor to the main refurbishment contractor. The scaffold had been erected without reference to TG20:21 and had not been formally inspected before use. The principal refurbishment contractor had not asked to see a scaffold inspection record, had not specified TG20:21 compliance in the scaffold subcontract, and had not verified the scaffolder's competence. The top lift of the scaffold was erected without a full set of standards and transoms, producing a configuration with excessive bay widths and no toe boards on the boarding. A window installer working from the top lift lost footing when reaching across an unbounded gap in the boarding and fell 7.2\,m to the pavement below, sustaining fatal head injuries. The investigation found that the principal contractor had not treated the scaffold as a safety-critical structure requiring design, inspection, and competence verification: it had been treated as an informal service procured at lowest cost with no governance. The working at height assessment existed as a generic document in the principal contractor's management system but had not been applied to the specific scaffold configuration, height, or task being undertaken.}

%---------------------------- CHAPTER SUMMARY ----------------------------
\newpage
\section{Chapter Summary}

\begin{table}[H]
\centering
\begin{tabularx}{\textwidth}{>{\raggedright\arraybackslash}p{3.2cm}X}
\toprule
\textbf{Assessment Element} & \textbf{Operational Requirement} \\
\midrule
Legal/Professional Basis & Work at Height Regulations 2005; NASC TG20:21 (tube-and-fitting scaffold); NASC SG4:22 (preventing falls in scaffolding operations); PASMA Code of Practice (mobile towers); BS EN~363 (personal fall protection systems); BS EN~795 (anchor devices); BS EN~1263 (safety nets); BS EN~131 (ladders); BS EN~361 (harnesses); BS EN~355 and EN~360 (lanyards and inertia reels). \\
Method & Apply the full five-step process and document inherent/residual risk. Assessments must be task-specific and trade-specific --- a generic working at height assessment does not meet the regulatory standard. The WAHR~2005 hierarchy (avoid, prevent, mitigate) must structure the control selection process. \\
Controls & Collective protection preferred over individual: TG20:21-compliant scaffold with 950\,mm top guardrail, intermediate rail and 150\,mm toe board; safety netting to BS EN~1263; proprietary safety decking during concrete pours; PASMA-compliant mobile towers; MEWP from IPAF-registered contractor; BS EN~363 personal fall protection systems where collective measures are impracticable; seven-day scaffold inspection with signed records; SG4:22 advanced guardrail sequence for scaffold erection and dismantling. \\
Cadence & Seven-day scaffold inspection mandatory; inspect before first use after erection or modification; inspect after adverse weather; review assessment at each new trade mobilisation, phase change, or programme sequence change; review immediately following any fall or near miss. \\
Assurance & Use audit, incident learning, and governance reporting to verify effectiveness. Maintain scaffold inspection register (retained for three months post-dismantling); MEWP daily pre-use log; personal fall protection equipment individual log; PASMA training records; competence evidence for scaffold inspector as part of the project safety file. \\
\bottomrule
\end{tabularx}
\end{table}

\begin{itemize}
  \bitem{Key takeaway 1}{The Work at Height Regulations 2005 apply to all work at height where there is a risk of a fall liable to cause personal injury, with no minimum height threshold. Fatal falls on construction sites occur from heights as low as 2\,metres. The assessment must address all height work regardless of duration or perceived triviality.}

  \bitem{Key takeaway 2}{Collective protective measures --- edge protection, safety netting, scaffold working platforms --- must be selected in preference to personal fall protection systems wherever they are technically and practically achievable. Personal fall protection is the last line of defence, not the default solution, and its effectiveness depends entirely upon correct selection, fitting, connection to a suitable anchor, and maintenance.}

  \bitem{Key takeaway 3}{Scaffold must be designed to TG20:21, erected and dismantled following the SG4:22 advanced guardrail sequence, formally inspected by a competent person before first use and at seven-day intervals thereafter, and have its inspection records retained for three months after dismantling. The scaffold is a safety-critical structure and must be governed as such, including through the subcontractor selection and supply chain competence verification processes.}

  \bitem{Key takeaway 4}{MEWPs used as a working at height access method introduce the additional hazards of lifting operations and must be managed in accordance with the guidance in Chapter~\ref{chap:Lifting Operations and Lifting Equipment Risk Assessment} as well as the working at height provisions of WAHR~2005. MEWP operators must hold the relevant IPAF card and the machine must undergo daily pre-use inspection and formal periodic thorough examination under LOLER~1998.}

  \bitem{Key takeaway 5}{Fragile surfaces --- rooflights, profiled metal decking prior to concrete topping, fragile roof coverings, and partially completed floor slabs --- present a specific and frequently fatal working at height hazard that must be explicitly identified and addressed in the risk assessment. The detailed treatment of roofing and fragile surface risks, including appropriate collective and individual protection measures, is provided in Chapter~\ref{chap:Roofing and Fragile Surfaces Risk Assessment}.}
\end{itemize}

\barquote{In Chapter~\ref{chap:Confined Spaces Risk Assessment}, the assessment methodology is applied to confined spaces --- an environment that reverses the spatial character of height work but shares its defining characteristic: the worker's ability to escape quickly from a developing emergency is severely compromised, making prevention through rigorous pre-entry assessment and permit-to-enter control an absolute imperative.}
